\documentclass[11pt,a4paper]{article}
\usepackage[margin=25mm]{geometry}
\usepackage{amsmath,amssymb,amsthm,hyperref}
\newtheorem{lemma}{Lemma}
\newtheorem{theorem}[lemma]{Theorem}
\newcommand{\gen}[1]{\langle#1\rangle}
\newcommand{\Sym}{\operatorname{Sym}}
\title{Constructing the rank interval with sign changes}
\author{Human Compression project: F006 recovery}
\date{Author working proof; distinct review pending}
\begin{document}\maketitle
Use $E_n$ and $W_n=E_n\rtimes S_n$ from the companion note
\emph{Even sign changes and the rank of a string}. The permutation
subgroup $S_n$ acts without sign changes. The vector $e_i+e_j$ flips
the two signed coordinates $i,j$. We use the standard end criterion:
a string of involutions is a C-marking if its two end deletions are
C-markings and their groups intersect in the common middle group.
This is the retained McMullen--Schulte criterion, Proposition~2E16.

\begin{lemma}[Adjoining a kernel involution]\label{kernel}
Let $(s_1,\ldots,s_d)$ be a string C-marking of a complement $S$ in
a split extension $N\rtimes S$. Let $x\in N$ be a nonidentity
involution commuting with $s_1,\ldots,s_{d-1}$. Then
$(s_1,\ldots,s_d,x)$ is a string C-marking of its generated subgroup.
\end{lemma}
\begin{proof}
Projection to $S$ shows that, for every $k$,
\[
 S\cap\gen{s_k,\ldots,s_d,x}=\gen{s_k,\ldots,s_d}.
\]
Indeed the projected image is the group on the right, and projection
is the identity on $S$; the reverse containment is explicit.
The pair $(s_d,x)$ consists of distinct nonidentity involutions and
has the intersection property. Induct backwards on $k$. One end
of $(s_k,\ldots,s_d,x)$ is an old C-submarking in $S$; the other
is the new suffix already proved C. The displayed identity makes
their intersection precisely $\gen{s_{k+1},\ldots,s_d}$.
The end criterion proves the induction. The case $d=1$ is the
initial pair. Generation of the whole extension is a separate question.
\end{proof}
This proves the intersection step of Chen--Rowley, Lemma~18, for an
arbitrary split kernel. No assertion about intersections of projected
subgroups is being substituted for the actual intersection property.

\section{A symmetric marking with an alternating path at its end}
For $n\ge5$ and $3\le d\le n-2$, put
\[
 a_i=(i\ i+1)\ (1\le i\le d-2),\qquad
 b=\prod_{j=d-1,d+1,\ldots<n}(j\ j+1),\qquad
 c=\prod_{j=d,d+2,\ldots<n}(j\ j+1).
\]
\begin{lemma}\label{symmetric}
The tuple $(a_1,\ldots,a_{d-2},b,c)$ is a string C-marking of $S_n$
with adjacent orders $[3^{d-3},6,n-d+2]$.
\end{lemma}
\begin{proof}
The union of the two matchings $b,c$ is the path on
$\{d-1,\ldots,n\}$. Its product is a cycle of length $n-d+2$,
as is seen by following alternate edges from an endpoint.
Thus $\gen{b,c}$ is transitive on that set. Conjugating the bridge
$(d-2\ d-1)$ by this group gives the transpositions joining $d-2$
to every point of the set, and hence its full symmetric group together
with $d-2$. The preceding bridges then give $S_n$.

For the intersection property first take $d=3$. Write $a=(1\ 2)$,
$b=(2\ 3)t$, where $t=(4\ 5)(6\ 7)\cdots$ is nontrivial and
disjoint from $\{1,2,3\}$. Since $(ab)^3=t$,
\[
 H=\gen{a,b}=S_{\{1,2,3\}}\times\gen t.
\]
The subgroup $K=\gen{b,c}$ is the dihedral action on the
$(n-1)$-cycle given by $bc$, and fixes $1$. Therefore
$H\cap K\le H_1=\gen{b,t}$. The points $2,3$ are adjacent in
that cycle and are both fixed by $t$. A nonidentity dihedral element
cannot fix two adjacent cycle points when the length is at least four:
a rotation fixes none, and the equations for a reflection fixing
successive residues would give $2=0$ modulo the cycle length.
Thus $t\notin K$ and $H\cap K=\gen b$. The end criterion gives
the rank-three intersection property.

For $d>3$, deleting $a_1$ gives the same construction on
$\{2,\ldots,n\}$ with rank $d-1$. Its first mark is the pure
transposition $(2\ 3)$, and every other mark fixes $2$.
Apply the acquired pure-leaf corollary of \emph{Restriction kernels
in extensions of string markings}, adjoining $(1\ 2)$ at a new point.
That corollary includes the complete intersection proof, with no
transitivity assumption on a projected action. Induction proves C.
Finally the products of consecutive pure transpositions have order
three; $a_{d-2}b$ is a three-cycle times a nonempty disjoint matching,
of order six; and the order of $bc$ was computed above. All other
required pairs commute directly from their supports.
\end{proof}
This is the explicit family attributed by Chen--Rowley, Theorem~17,
to Fernandes--Leemans (2011), Lemma~21. Its proof here uses the
already acquired restriction-kernel theorem, with the base intersection
checked directly.

\section{Odd degree}
\begin{theorem}\label{odd}
If $n\ge5$ is odd, then $W_n$ has a rank-$n$ C-marking of type
$[4,3^{n-2}]$ and C-markings in every rank from three through $n$.
\end{theorem}
\begin{proof}
Take $v=e_2+\cdots+e_n\in E_n$ and $s_i=(i\ i+1)$.
The adjacent transpositions form the usual C-marking of $S_n$:
every subset generates the independent symmetric groups on its path
components, whose intersections are the components of the intersection
of the two edge sets. The vector $v$ commutes with $s_i$ for $i\ge2$.
Lemma~\ref{kernel}, with the ordering reversed, proves C for
$(v,s_1,\ldots,s_{n-1})$.
Its group contains $S_n$ and the nonconstant vector $v$, so the
module lemma gives all of $E_n$. It is $W_n$.
The square of $v s_1$ is $v+s_1v=e_1+e_2\ne0$, so its order is
four; the other adjacent orders are three.

Apply the acquired Brooksbank--Leemans contraction theorem, in
\emph{Contraction of string markings through a central factor},
Theorem~3 and Corollary~5. Every bond tested in successive contractions,
starting at the third bond, has order three. Each contraction preserves
the group and the intersection property, and iteration ends at rank
three. This gives precisely the stated interval.
\end{proof}

\section{Even degree}
\begin{theorem}\label{even}
If $n\ge6$ is even, then $W_n$ has C-markings in every rank from
three through $n-1$. In rank three the type is $[12,n-1]$;
in rank $4\le r\le n-1$ it is $[3^{r-4},6,n-r+3,4]$.
\end{theorem}
\begin{proof}
For $r\ge4$ use Lemma~\ref{symmetric} with $d=r-1$.
Let $(u\ u+1)$ be the last transposition in $b$ and put
$v=e_u+e_{u+1}$. This is $(n-1\ n)$ when $r$ is odd and
$(n-2\ n-1)$ when $r$ is even. In either case $u>r-2$.
Thus $v$ commutes with all pure marks and with $b$.
Append $v$ after $c$ and apply Lemma~\ref{kernel}.
The group contains $S_n$ and a nonconstant vector, hence is $W_n$.
The matching $c$ does not preserve $\{u,u+1\}$, so
$(cv)^2=v+cv$ is a nonzero vector of order two. This gives the
last bond order four and proves the asserted type.

For rank three put
\[
 a=(e_{n-1}+e_n)(1\ 2),\quad
 b=(2\ 3)(4\ 5)\cdots(n-2\ n-1),\quad
 c=(3\ 4)(5\ 6)\cdots(n-1\ n).
\]
These are involutions. The supports of $(1\ 2)$ and $c$ are
disjoint, and $c$ preserves the sign-change pair $\{n-1,n\}$,
so $a,c$ commute. On signed letters $q=ab$ has a three-cycle
on $\{1,2,3\}$ and its negative copy, a four-cycle on
$\{\pm(n-2),\pm(n-1)\}$, a sign flip at $n$, and the
remaining disjoint unsigned transpositions. Hence $|q|=12$
and $q^6=e_{n-2}+e_{n-1}$. Also $bc$ is a cycle of length
$n-1$ on $\{2,\ldots,n\}$ and its negative copy.

The projected group contains $(1\ 2)$ and the transitive group
$\gen{b,c}$ on the other letters. Conjugates of $(1\ 2)$ give
all the star transpositions, so the image is $S_n$. The kernel
contains the nonconstant $q^6$. It is invariant under this full
image; by the module lemma it is $E_n$. Generation follows.

For the actual end intersection, every element of $A=\gen{a,b}$
has the unique form $b^\epsilon q^k$, $\epsilon\in\{0,1\}$,
$0\le k<12$. Since $b$ is unsigned, such an element is unsigned
exactly when $q^k$ is unsigned, which the signed four-cycle shows
to be equivalent to $4\mid k$. Since $b$ fixes $1$, the same
element fixes $1$ exactly when $3\mid k$, by the three-cycle.
Therefore
\[
 A\cap S_n\cap\operatorname{Stab}(1)=\gen b.
\]
As $\gen{b,c}$ is contained in the middle group on the left,
$\gen{a,b}\cap\gen{b,c}=\gen b$. Both end pairs are C,
so the end criterion proves the full rank-three assertion.
\end{proof}

The last argument corrects the displayed subgroup calculation in
Chen--Rowley, Lemma~19. Both subgroups in its displayed last
intersection contain $t_1$, although $t_1$ is not unsigned and the
claimed intersection is $\gen{t_2}$. The actual intersection needed
by the lemma is indeed $\gen{t_2}$, as the explicit calculation above
proves. The rank-three theorem itself is recovered.

Combining Theorems~\ref{odd} and~\ref{even} with the companion upper
bound proves the full principal $D_n$ rank conclusions of source
Theorems~1 and~2. All constructions, types and the contraction method
remain attributed to their original authors. The reorganization uses
one kernel-extension lemma and the acquired module, restriction-kernel
and contraction arguments. Full dependency proofs and their retained
inputs remain required. The exceptional census and distinct review
are not completed by this working note; no whole-manuscript or numerical
compression claim is attached to it.
\end{document}
